
\lussec 3. $\mib D$+1 维场论

以下各节的讨论在很大程度上基于这样一个事实：
随时间演化场的关联函数与
$D$+1 维局域场论中的关联函数一致，
多出的一维即流时间。
这一观察可追溯到 Zinn--Justin 与 Zwanziger 关于
朗之万方程的开创性工作 [\cite{Zinn},\cite{ZinnZwanziger}]，
并且最近使得（纯规范的）梯度流的重正化性
能够在微扰论的所有阶得以建立
[\cite{RenFlow}]。

如下文将变得清楚的，
在 $D$+1 维理论中纳入夸克场是直截了当的。
总体框架与文献~[\cite{RenFlow}] 完全相同，
那里可以找到对该专题的导引及更多细节。


\lussubsec 3.1 场与作用量

除已经遇到的基本场和随时间演化的场之外，
$D$+1 维理论还涉及拉格朗日乘子场
$L_{\mu}(t,x)$、$\lambda(t,x)$ 与 $\lambdabar(t,x)$。
后者是与夸克场带有相同指标的费米子场，
而 $iL_{\mu}(t,x)$ 是取值于 $\SUn$ 李代数的矢量场。

如文献~[\cite{RenFlow}] 所解释的，
微扰论中的规范固定要求引入法捷耶夫--波波夫鬼场
$c(x),\cbar(x)$ 及相应的随时间演化场 $d(t,x),\dbar(t,x)$，其中
\equation{
  \left.d\right|_{t=0}=c.
  \enum
}
场 $\dbar(t,x)$ 不施加边界条件；
它在许多方面起着类似于拉格朗日乘子场的
作用。
除边界条件 (2.1)、(2.6) 与 (3.1) 外，
此处所有场均被假定为不受约束的。

理论的 total 作用量
\equation{
  S_{\rm tot}=S+\SGfl+\SFPfl+\SFfl,
  \enum
}
包含流时间为零的场的规范固定 QCD 作用量 $S$
以及体作用量
\equation{
  \SGfl=-2\int_0^{\infty}\rmd t\int\rmd^D\kern-1pt x\, 
  \tr\bigl\{L_{\mu}(t,x)
  \bigl(\partial_tB_{\mu}-
        D_{\nu}G_{\nu\mu}-\alpha_0D_{\mu}\partial_{\nu}B_{\nu}
  \bigr)(t,x)\bigr\},
  \enum
  \nexteq{3.0ex}
  \SFPfl=-2\int_{0}^{\infty}\rmd t\int\rmd^D\kern-1pt x\,
  \tr\bigl\{
  \dbar(t,x)\bigl(\partial_td-\alpha_0 D_{\mu}\partial_{\mu}d\bigr)(t,x)
  \bigr\},
  \enum
  \nexteq{3.0ex}    
  \SFfl=\int_0^{\infty}\rmd t\int\rmd^D\kern-1pt x\,
  \bigl\{\lambdabar(t,x)
         (\partial_t-\Delta+\alpha_0\partial_{\nu}B_{\nu})\chi(t,x)
  \noenum
  \nexteq{1.5ex}
  {\phantom{\SFfl=\int_0^{\infty}\rmd t\int\rmd^D\kern-1pt x\,\bigl\{}}
  +\chibar(t,x)
  (\lvec{\partial}_t-\lvec{\Delta}-\alpha_0\partial_{\nu}B_{\nu})
  \lambda(t,x)\bigr\}.
  \enum
}
相对于文献~[\cite{RenFlow}] 中研究的纯规范流，
体夸克作用量 $\SFfl$ 是 total 作用量中唯一的新增项。


\lussubsec 3.2 关联函数

规范场、夸克场、鬼场与拉格朗日乘子场的
$n$ 点关联函数，
形式上通过带``玻尔兹曼因子''
$\exp\{-S_{\rm tot}\}$ 的对全体场的泛函积分来定义。

对于规范场与夸克场的关联函数，
泛函积分可以在一定程度上通过先对拉格朗日乘子场
积分而算出。
由于作用量对这些场是线性的，
积分给出泛函狄拉克 $\delta$ 函数的乘积。
体作用量的选取使得这些 $\delta$ 函数
等价于把流方程强加于随时间演化的规范场与夸克场。
除数学上的微妙之处以及可能非平凡的雅可比因子外，
这样得到的关联函数与第 2 节中直接计算的结果一致：
即用初始条件 (2.1)、(2.6) 求解流方程，
并把算出的随时间演化场当作 QCD 可观测量。

$D$+1 维理论与直接计算关联函数之间的等价性，
可在微扰论中 [\cite{RenFlow}]
及格点 QCD 中（第 5 节）严格证明。
特别地，所有不涉及随时间演化夸克场
或相应拉格朗日乘子场的关联函数，
都与文献~[\cite{RenFlow}] 中研究的（夸克场不随流时间演化的）
关联函数一致。
$D$+1 维理论的这一``约化性质''
强烈限制了重正化所需的可能抵消项的形式
（见 3.4 小节）。


\lussubsec 3.3 费米子积分

由于 $D$+1 维中的作用量对费米子场是二次的，
这些场的泛函积分可以通过应用 Wick 定理求值（固定玻色场）。
基本 Wick 收缩值得明确写出，
因为这些表达式提供了对 $D$+1 维理论结构的洞察，
而且无论如何以后会经常用到。

基本场的 Wick 收缩由作用量和边界条件 (2.6) 决定。
特别地，流时间为零处夸克场的传播子
\equation{
  \wick{\psi(x)\psibar(y)}=S(x,y),
  \enum
  \nexteq{2.5ex}
  (\slash{D}+M_0)S(x,y)=\delta(x-y),
  \enum
}
与 $D$ 维中的传播子一致。
除夸克传播子 $S(x,y)$ 外，
随时间演化费米子场的收缩
\equation{
  \longwick{\lambda(t,x)\lambdabar(s,y)}=0,
  \enum
  \nexteq{2.5ex}
  \longwick{\chi(t,x)\lambdabar(s,y)}=\theta(t-s)K(t,x;s,y),
  \enum
  \nexteq{2.5ex}
  \longwick{\lambda(t,x)\chibar(s,y)}=\theta(s-t)K(s,y;t,x)^{\dagger},
  \enum
  \nexteq{2.5ex}
  \longwick{\chi(t,x)\chibar(s,y)}=\int\rmd^Dv\,
  \rmd^Dw\,
  K(t,x;0,v)S(v,w)K(s,y;0,w)^{\dagger},
  \enum
}
涉及夸克流方程的基本解
\equation{
  \{\partial_t-\Delta+
  \alpha_0\partial_{\nu}B_{\nu}\}K(t,x;s,y)=0\quad\hbox{if}\quad t\geq s,
  \enum
  \nexteq{2.0ex}
  \lim_{t\to s}K(t,x;s,y)=\delta(x-y),
  \enum
}
注意核 $K(t,x;s,y)$ 与味道无关，
正比于单位狄拉克矩阵，
且是颜色空间中的 $N\times N$ 复矩阵
（式 (3.10)、(3.11) 中核的伴随仅在指标空间上取）。

与流时间为零处的边界条件一致，
涉及基本夸克场与反夸克场的收缩
\equation{
  \longwick{\chi(t,x)\psibar(y)}=\int\rmd^Dv\,K(t,x;0,v)
  S(v,y),
  \enum
  \nexteq{2.5ex}
  \wick{\psi(x)\chibar(s,y)}=\int\rmd^Dw\,S(x,w)
  K(s,y;0,w)^{\dagger},
  \enum
  \nexteq{2.5ex}
  \longwick{\lambda(t,x)\psibar(y)}=
  \wick{\psi(x)\lambdabar(s,y)}=0,
  \enum
}
作为特例已包含在式 (3.9)--(3.11) 中。
然而要注意，
时空边界附近关联函数的行为
往往会因短距离奇异性而复杂化
[\cite{Symanzik}]。
例如极限
\equation{
  \lim_{s\to0}\longwick{\lambda(0,x)\chibar(s,y)}=
  \lim_{t\to0}\longwick{\chi(t,x)\lambdabar(0,y)}=
  \delta(x-y),
  \enum
}
在 $t=s=0$ 处与收缩 (3.16) 相差一个接触项。


\lussubsec 3.4 重正化

在微扰论中，
$D$+1 维理论的重正化性质可以用幂计数
与对称性论证来确定。
由于已知只让规范场随流时间演化的理论，
其重正化不需要超出通常的 QCD 参数与场重正化
[\cite{RenFlow}]，
找出作用量全部可能的局域抵消项的工作因此得到简化。
特别地，在流时间 $t>0$ 处，
场 $B_{\mu}(t,x)$ 与 $L_{\mu}(t,x)$ 无需重正化。

当随时间演化的夸克场被纳入理论时，
可能需要向 total 作用量添加进一步的抵消项。
然而 3.2 小节末提到的约化性质意味着，
任何这样的项都必须正比于夸克场或
相应的拉格朗日乘子场。
此外如文献~[\cite{RenFlow}] 的 7.2 小节所解释的，
体抵消项（即在全部 $D$+1 个坐标上积分的局域项）是被排除的，
因为大流时间下的关联函数由树图给出。

考虑到未重正化理论的对称性，
以及夸克场与拉格朗日乘子场的幂计数量纲分别为 $3/2$ 和 $5/2$，
稍加思考即可发现唯一可能新增的抵消项正比于
\equation{
  \int\rmd^D\kern-1pt x\,
  \left\{\lambdabar(0,x)\psi(x)+\psibar(x)\lambda(0,x)\right\}.
  \enum
}
该项在费曼图中的插入，
等价于正流时间下费米子场
以及类似地反费米子场的一个重正化
\equation{
  \chi=Z_{\chi}^{-1/2}\ren{\chi},
  \qquad
  \lambda=Z_{\chi}^{1/2}\ren{\lambda},
  \enum
}
因此，除 QCD 参数与场的通常重正化之外，
$D$+1 维理论的重正化还要求按上述方式
重正化随时间演化的夸克场及相应的拉格朗日乘子场。
